\section{Implementation specification}\label{app:impl}

\begin{revblock}
This appendix specifies the deployed procedure completely, so that it
can be re-implemented from the paper alone; the released code
(Section~``Code and data availability'') is the executable reference and
fixes every random seed used in the experiments.

\paragraph{Inputs and preprocessing.}
The method takes a feature image $\bm z$ on a lattice $D$ (any $p$); no
preprocessing is part of the method itself. In all real-data experiments the
bands are standardised to zero mean and unit variance per band, and the
hyperspectral cubes are reduced by principal components to $p=10$
(Section~\ref{subsec:data}); Section~\ref{subsec:realspec} reports the
insensitivity of the results to that truncation ($5$ to $30$ components).

\paragraph{Fixed constants.}
All constants are fixed once, across every experiment in the paper; none is
tuned per scene.
\begin{center}
\small
\begin{tabular}{llp{4.2cm}l}
\toprule
Symbol & Value & Role & Where it enters\\
\midrule
$\alpha$ & $0.05$ & test level & Eq.~\eqref{eq:test}\\
$\tau$   & $0.25$ & detrend threshold & Eq.~\eqref{eq:detrendrule}\\
--       & $2$    & polynomial degree of the drift & Eq.~\eqref{eq:r2}\\
$t$      & $24$   & tile side in pixels (reduced to
                    $\min(24,\lfloor H/2\rfloor,\lfloor W/2\rfloor)$ on
                    small domains) & Eq.~\eqref{eq:areamed}\\
$\nu$    & $12$   & minimum effective size to test a node &
                    Algorithm~\ref{alg:essdip}\\
$B$      & $5$    & ensemble size (median of $\hat K$) &
                    Section~\ref{subsec:algo}\\
--       & $300$  & Monte Carlo draws per null quantile &
                    Section~\ref{subsec:ess}\\
--       & $[12,2000]$ & clamp of the rounded $\neff$ before quantile lookup &
                    Section~\ref{subsec:ess}\\
--       & $8$ / $20$ & recursion cap on $\hat K$ (factorial study / full
                    scenes) & Algorithm~\ref{alg:essdip}\\
\bottomrule
\end{tabular}
\end{center}

\paragraph{Null quantile.}
For a node with effective size $\neff$, let
$m=\operatorname{clamp}(\lfloor\neff\rceil,12,2000)$. The critical value
$q_{1-\alpha}(m)$ is the empirical $(1-\alpha)$ quantile of $300$ dip
statistics, each computed from $m$ pseudo-random draws of the reference null:
$\mathrm{U}(0,1)$ for ESS-Dip (least-favourable, distribution-free) and
$\mathrm{N}(0,1)$ for ESS-Dip-R (the null of Assumption~\ref{ass:null}).
Quantiles are memoised by $m$ (and by reference law), so each distinct size is
simulated once per run.

\paragraph{Split proposal.}
At a node $C$, $2$-means is run on the feature vectors of $C$ ($k$-means with
$k=2$, five restarts); the split direction $\bm u$ is the normalised
difference of the two cluster means, Eq.~\eqref{eq:proj}, and the dip is
computed on $\{\bm u^{\!\top}\bm z(s)\}_{s\in C}$ by the algorithm of
\citet{HartiganHartigan1985}. A significant test splits $C$ into the two
$2$-means children, each retained only if its pixel count is at least the
decorrelation area $a$ (so that $\neff\ge1$).

\paragraph{Range estimation.}
The decorrelation area $a=\ell^2$ is estimated once per scene (not per node),
on the detrended field when Eq.~\eqref{eq:detrendrule} activates: per tile,
the first principal component's sample autocorrelation is computed by FFT and
its $1/e$ crossing recorded, averaged over the two axes; $\ell$ is the median
over tiles, Eq.~\eqref{eq:areamed}.

\paragraph{Ensemble.}
The recursion is run $B=5$ times with independent seeds (the only stochastic
elements are the $k$-means initialisations and the null-quantile Monte Carlo);
the reported $\hat K$ is the ensemble median.
\end{revblock}
